\documentclass{situationist-preprint}

\title{Exceptions and Generalizations}
\msid{SCI-08-b}
\msfamily{SCI-08}
\msdiscipline{Philosophy of Science}
\mstype{Companion manuscript}
\mskeywords{biological generalizations; exceptions; ceteris paribus; meiotic drive; laws of nature}

\begin{document}
\maketitle

\begin{abstract}
SCI-08-a argues that there are no strict laws in biology, but robust generalizations with stable exceptions. A standing objection to any generalization with exceptions is that it is either false or, once hedged by a ceteris paribus clause, empty. We argue that the objection fails for generalizations whose exceptions are stable: because the exceptions fall under projectible generalizations of their own, the generalization together with its exception-generalizations has definite content and can be tested. Mendel's law of segregation and the co-occurrence of the mammalian characters illustrate the two sources of stable exceptions in biology, deviant mechanisms and divergent lineages.
\end{abstract}

\section{The dilemma}\label{s:dilemma}

\citet{EarmanRoberts1999} pressed a dilemma on laws with exceptions. Stated without qualification, \enquote{all $F$ are $G$} is false if some $F$ are not $G$. Qualified as \enquote{ceteris paribus, all $F$ are $G$}, it is true only if the clause has content, and if the clause means \enquote{unless something interferes}, it is true whatever happens. Either the generalization is false or it is empty.

SCI-08-a accepts the first horn for biology: by its Theorem~3.1 there are no strict biological laws. Its Proposition~3.2 claims that some biological generalizations are robust and have only stable exceptions, in the sense of Definitions~2.2 and~2.3 of SCI-08-a. We show that such a generalization escapes the second horn.

\begin{proposition}\label{p:content}
If the exceptions to a generalization $G$ are stable, then $G$ together with the generalizations covering its exceptions entails observational consequences that would be false if $G$ failed in cases not covered by them.
\end{proposition}

\begin{proof}
By Definition~2.3 of SCI-08-a, a stable exception is produced by an identifiable and recurrent mechanism, so the class of exceptions is the subject of projectible generalizations. The conjunction of $G$ with these generalizations says that every $F$ is $G$ unless one of the identified mechanisms is present, and that where it is present the outcome is as the exception-generalization says. The conjunction is false of any $F$ that is not $G$ and in which no identified mechanism is present. Unlike \enquote{unless something interferes}, the qualification names its interferers in advance, and so has content.
\end{proof}

\section{Two sources of stable exceptions}\label{s:sources}

\paragraph{Deviant mechanisms.} Mendel's law of segregation says that the two alleles at a locus of a heterozygote are transmitted to gametes in equal proportions. It fails in meiotic drive, in which a segregation distorter, such as the $t$-haplotype in the house mouse or the \emph{Segregation Distorter} system in \emph{Drosophila}, is transmitted to most of the functional gametes \citep{Lyttle1991}. The exceptions are stable: the mechanisms are identified, they recur across taxa, and their effects are the subject of their own generalizations. The law of segregation together with the generalizations of drive predicts that a locus with no distorter will show equal transmission, and a locus that departs from it is evidence of a distorter to be found.

\paragraph{Divergent lineages.} The characters of mammals, hair, mammary glands, three middle-ear bones, co-occur across the class, and the generalization that mammals bear live young holds with the exception of the monotremes. The co-occurrence is explained by common descent: the characters were inherited together from a common ancestor and have been retained in most lineages, and the exceptions are lineages that diverged before a character arose or that lost it. This generalization is robust and its exceptions are stable, since they fall under the generalizations of phylogeny. But it is not nomic. There are mammals without the character, the monotremes having diverged before it arose, and the co-occurrence of the characters holds because of shared history, not because any of them requires the others. MET-27-b draws the consequence for natural kinds: Mammalia passes the outsider test, but it is probably not the extension of a nomic cluster.

\begin{objection}[Strict laws in disguise]
If the exception-generalizations are added to the original, the conjunction is an exceptionless generalization in biological vocabulary, and so a strict law.
\end{objection}

\begin{reply}
The conjunction is exceptionless only relative to the mechanisms identified so far, and by Assumption~2.5 of SCI-08-a variation can produce new ones, so that its own exceptions are not excluded. What is stable is the pattern of exceptions known at a time, not a final list. The conjunction is a better generalization, not a strict law, which is the distinction between robustness and strictness that SCI-08 draws.
\end{reply}

\section{Conclusion}\label{s:conclusion}

The dilemma of falsity or vacuity applies to generalizations hedged by an unspecified ceteris paribus clause. It does not apply to biological generalizations whose exceptions are stable, because those exceptions are named in advance by mechanism or by lineage. Biology has no strict laws, but robust generalizations with stable exceptions, as SCI-08 concludes, and the exceptions are part of what gives the generalizations their content.

\begin{thebibliography}{9}
\bibitem[Earman and Roberts(1999)]{EarmanRoberts1999} Earman, J., and Roberts, J. (1999). Ceteris paribus, there is no problem of provisos. \emph{Synthese}, 118, 439--478.
\bibitem[Lyttle(1991)]{Lyttle1991} Lyttle, T.~W. (1991). Segregation distorters. \emph{Annual Review of Genetics}, 25, 511--557.
\end{thebibliography}

\end{document}
